\section{Materials and Methods}

\subsection{Study design}

We designed the study as a representation-diagnostic analysis of controlled short-read perturbation settings. Here, representation diagnostics denotes a prespecified set of representation-level measurements rather than a clinical decision rule: paired stability, nearest-clean retrieval, grouped local-change readout, feature dimension and blockwise decomposition. The central question was whether biochemical and coarse positional summaries add auditable evidence to a canonical local k-mer-composition backbone. CK4 and CK5 were reverse-complement canonical k-mer-composition baselines; P denoted global biochemical-property summaries; MSP denoted multi-scale positional property pooling. The seven non-empty K/P/MSP combinations were evaluated, full-position matrices were treated as positional upper bounds, and the canonical spaced-property (CSP) control served as a spaced-seed mechanism comparator.

This design makes two questions testable. First, whether a compact mixed representation remains close to its clean counterpart under nuisance perturbation while retaining representation-level readability. Second, whether the observed behavior can be assigned to metric-specific K, P or MSP contributions rather than to dimensionality, normalization or dilution in a mixed L2 space. We therefore combined a seven-group ablation with same-dimension principal-component analysis (PCA) and singular-value decomposition (SVD) controls, dimension-matched high-k compressed k-mer baselines, block-weight sensitivity, paired non-parametric tests, k-nearest-neighbour mutual-information robustness checks, P/MSP relation analysis, feature-extraction runtime and error-aware ART simulator strata. Shallow readout probes tested whether signal remained linearly accessible; they were not treated as production classifier performance. Throughout, mixed-space L2 is interpreted as standardized diagnostic drift after block construction and normalization, not as a natural biophysical distance between commensurate physical units.

The 69 and 75 bp settings fall within a commonly reported 50--75 bp mNGS read-length range \citep[p.~99]{Yang2026mNGS} and correspond to aggregate pre- and post-quality-control mean-length conditions observed in a restricted local sequencing context. Only these aggregate length conditions informed the design; no patient-level reads, labels, identifiers or sequence content were analysed. The two lengths were used as lower-read-length stress conditions, not as a claim that short reads are unusable. The remaining lengths were selected for sensitivity analysis: 100, 125 and 150 bp test whether the trends persist as reads become less extreme, and 300 bp is used only as an extended-read reference in selected boundary analyses.

\begin{figure*}[!tbp]
\centering
\includegraphics[width=0.92\textwidth]{figures/main/nature_fig1_framework.pdf}
\caption{Representation-diagnostic framework and read-length regime. Clean and perturbed reads are encoded through canonical local k-mer composition (K), global biochemical summaries (P) and multi-scale positional property pooling (MSP). The audit reports paired stability, nearest-clean retrieval, grouped local-change readout and feature dimension across 69--150 bp conditions; full-position encodings and the canonical spaced-property (CSP) control provide boundary comparisons.}
\label{fig:figure1}
{\small\noindent\textbf{Alt text:} Workflow diagram showing short reads separated into local-composition, biochemical and positional representation channels before diagnostic readouts and boundary checks.\par}
\end{figure*}


\subsection{Representation families}

The representation families were defined to separate local k-mer composition, global biochemical summaries and coarse positional information. CK4 and CK5 count reverse-complement canonical contiguous k-mers. P summarizes hydrogen-bond class, GC content, purine content, electron-ion interaction potential (EIIP)-like values, N fraction, entropy and length-related properties at the read level. MSP applies per-base property channels over relative-position bins. Its purpose is not to create a larger sequence vocabulary or to reconstruct full position. Instead, it adds a coarse positional property summary to an otherwise order-poor k-mer-frequency vector. The default 2+3+4+6 binset was prespecified as a coarse-to-fine relative-position summary whose finest scale remains above single-position resolution across the 69--150 bp range. The seven-group ablation comprised K (136 dimensions), P (11), MSP (75), K+P (147), K+MSP (211), P+MSP (86) and K+P+MSP (222). Full-position composition or property matrices retained per-position information and were used only to estimate an upper bound on positional readability. CSP combines spaced-seed counts with property summaries and was used to test whether matching-oriented spaced-seed priors transfer into dense representation diagnostics.

For mixed representations, each block was computed separately and internally L2-normalized before weighted assembly. For read $s_i$, let $\mathbf{K}_i$, $\mathbf{P}_i$ and $\mathbf{M}_i$ denote the reverse-complement canonical 4-mer composition block, the global biochemical-property block, and the multi-scale positional-property block, respectively. The internally normalized blocks were defined as

\[
\hat{\mathbf{K}}_i=\frac{\mathbf{K}_i}{\left\|\mathbf{K}_i\right\|_2},\quad
\hat{\mathbf{P}}_i=\frac{\mathbf{P}_i}{\left\|\mathbf{P}_i\right\|_2},\quad
\hat{\mathbf{M}}_i=\frac{\mathbf{M}_i}{\left\|\mathbf{M}_i\right\|_2},
\]

with zero-norm safeguards applied before normalization. The assembled compact representation was then

\[
\mathbf{x}_i
=
\frac{1}{\sqrt{\alpha^2+\beta^2+\gamma^2}}
\left[\alpha \hat{\mathbf{K}}_i \,;\, \beta \hat{\mathbf{P}}_i \,;\, \gamma \hat{\mathbf{M}}_i\right],
\]

where $\alpha$, $\beta$ and $\gamma$ are block weights and $[\cdot\,;\,\cdot]$ denotes concatenation. Unless otherwise stated, all three weights were set to 1. The block-weight audit compared equal weights with K-dominant $(2,1,1)$ and property-dominant $(1,2,2)$ settings; the separate MSP sensitivity audit evaluated $\gamma\in\{0,0.25,0.5,1,2\}$. Because the block norms are fixed before concatenation, the denominator is a fixed scalar for a given weight setting rather than a read-specific global norm. The reported global drift metric was

\[
d(i,j)=\left\|\mathbf{x}_i-\mathbf{x}_j\right\|_2,
\]

which should be interpreted as a standardized diagnostic drift in the assembled feature space, not as a natural physical distance between commensurate biological units. This definition also gives the block-normalized distance identity

\[
\begin{aligned}
d_{\mathrm{mix}}^2(i,j)
&=
\frac{1}{\alpha^2+\beta^2+\gamma^2}
\Bigl[
\alpha^2\left\|\hat{\mathbf{K}}_i-\hat{\mathbf{K}}_j\right\|_2^2 \\
&\quad+
\beta^2\left\|\hat{\mathbf{P}}_i-\hat{\mathbf{P}}_j\right\|_2^2 \\
&\quad+
\gamma^2\left\|\hat{\mathbf{M}}_i-\hat{\mathbf{M}}_j\right\|_2^2
\Bigr].
\end{aligned}
\]

For comparison with the composition-only drift $d_K(i,j)=\|\hat{\mathbf{K}}_i-\hat{\mathbf{K}}_j\|_2$, the mixed squared drift is lower than $d_K^2$ when

\[
\beta^2\left(d_P^2-d_K^2\right)+\gamma^2\left(d_M^2-d_K^2\right)<0,
\]

where $d_P$ and $d_M$ are the corresponding blockwise drifts. Thus, lower property and MSP drift than K drift is a sufficient, but not necessary, condition for lower mixed drift. The diagnostic protocol takes matched clean/perturbed reads, constructs each block without fitting a predictive encoder, reports paired stability and nearest-clean retrieval, tests structured local change with grouped delta-readout, and decomposes the result through the seven-group ablation and prespecified conditional contrasts. This protocol yields a representation profile rather than an automated acceptance threshold. The design therefore does not assume that concatenation is intrinsically useful; it requires the combination to survive dimensionality, weight, counterfactual and block-contribution audits.

\begin{table*}[!tbp]
\centering
\scriptsize
\caption{Representation families and diagnostic roles.}
\label{tab:table1}
\begin{tabular}{@{}p{0.18\textwidth}p{0.24\textwidth}p{0.24\textwidth}p{0.26\textwidth}@{}}
\toprule
\textbf{family} & \textbf{information channel} & \textbf{main role} & \textbf{main-text interpretation} \\
\midrule
CK4 / CK5 & canonical local k-mer composition & composition backbone and comparator & compact token-frequency evidence, not database-linked read identity \\
P & global property summary block & perturbation-stable channel & tests a global nucleotide-property and auxiliary summary without local composition \\
MSP & positionalized nucleotide-property summaries & coarse positional channel & tests grouped local-change readout without K or global P \\
CK4+P & composition plus global P summary & global property baseline & tests whether the complete global summary adds to CK4 \\
CK4+MSP & composition plus positionalized nucleotide-property signal & positional property ablation & tests whether MSP returns coarse layout information without global P \\
P+MSP & global plus positionalized property summaries & property-only combination & tests property stability and retrieval without K \\
CK4P-MSP & composition plus global and multi-scale property summaries & compact main representation & balances stability, grouped local-change readout and composition-linked retrieval \\
PseKNC / NCP+ANF / PseEIIP & historical composition--property, order or position descriptors & direct handcrafted-feature controls & define whether CK4P-MSP offers a distinct attributable stability--readout trade-off beyond conventional joint feature evaluation \\
full-position matrix & per-position token or property signal & high-dimensional positional-readout reference & tests fine positional information at higher dimensional cost \\
CSP & spaced-seed prior plus biochemical summary & spaced-seed transfer comparator & delimits where matching-oriented spaced seeds transfer to dense features \\
\bottomrule
\end{tabular}

\end{table*}


\subsection{Data layers and perturbation design}

The experiments were organized into eight data layers, with the scale of each layer reported in Table 2. The WGS perturbation layer contained 25,200 rows from 3,600 clean templates across six genera, 21 species and six read lengths. The position-property ablation layer expanded this to 36,000 rows over the same template set and length grid. ART Illumina-like simulation contributed 75,592 paired rows from 37,796 templates across 69--150 bp reads \citep{Huang2012ART}. Critical Assessment of Metagenome Interpretation (CAMI) TOY low-complexity probes used a 104,166-row initial labelled subset and a 216,000-row expanded subset with 30 candidate taxon identifiers at 69, 75 and 100 bp \citep{Sczyrba2017}. The CAMI II marine lightweight probe used 4,000 anonymous source reads and 48,000 length/condition rows at 69, 75 and 100 bp \citep{Meyer2022}. The local mutation layer used 250 triplets per analysis cell, four local modes and three read lengths. WGS-derived paired perturbations measured clean-versus-perturbed stability; ART added simulator-derived errors; CAMI resources supplied bounded external probes; controlled position and order tasks tested positional readability; and boundary probes tested spaced-seed transfer, context visibility and ARG/SNP limits.

For CAMI\_TOY\_low, read identifiers were joined to the available taxon mapping and normalized to a common label field; the expanded source subset contained 30 candidate taxon identifiers. Unclassified reads were excluded from the labelled probe and retained only when a target-versus-background split explicitly required background reads. Per-class sampling was capped to prevent dominant labels from controlling the shallow probe. These CAMI readouts were generated before the public block-normalized CK4P-MSP contract and are retained only as external context; historical feature aliases were not relabelled as the current method. Consequently, the main text uses the coarse target-versus-background result and does not present the historical 30-label probe as a contract-v2 CK4P-MSP comparison.

We also added a lightweight CAMI II marine subset as an external short-read stability probe. We streamed a prefix of CAMI II marine short-read sample 0 and parsed 4,000 anonymous 150 bp reads from the embedded \texttt{anonymous\_reads.fq.gz} member. These reads were truncated to 69, 75 and 100 bp and expanded into clean, \texttt{N\_3pct}, \texttt{substitution\_1pct} and \texttt{substitution\_1pct\_N\_3pct} conditions, yielding 48,000 length/condition rows. CAMI II provides benchmark truth resources, but the anonymous FASTQ headers did not provide read-level taxonomic labels in the streamed-read file used here, and the separate OTU-specific truth bundles were not reconstructed for this lightweight analysis. This CAMI II probe therefore reports paired perturbation-stability metrics on an external metagenomic sequence source rather than CAMI II taxonomic validation.

This layered simulation design was used because each component isolates a different failure mode. The controlled WGS grid makes perturbation type, read length and representation family directly comparable. ART adds a recognized Illumina-like simulator layer, but it is retained as a consistency check rather than treated as ground-truth clinical sequencing physics. The quality-stratified ART audit adds a more explicit error-aware view of the same simulator-derived layer. It supports statements about robustness across quality strata, but it does not by itself establish validation on clinical sequencing data.

\begin{table*}[!tbp]
\centering
\scriptsize
\caption{Data layers, diagnostic questions, metrics and claim boundaries.}
\label{tab:table2}
\input{tables/main/table2.tex}
\end{table*}


\subsection{Metrics and statistical analysis}

Perturbation stability was quantified by paired cosine similarity, paired L2 drift and nearest-clean retrieval. Representation readability was quantified by macro-F1 and accuracy from shallow logistic-regression or nearest-centroid probes. These readout metrics were accessibility probes for low-dimensional features, not production classifier performance. Compactness was measured by feature dimension. Local mutation analysis used selective-sensitivity ratios and grouped delta-readout; all derivatives of one template were kept in the same fold through stratified grouped cross-validation. For the seven-group audit, three contrasts were prespecified: K+P+MSP versus P+MSP for the conditional K contribution, versus K+MSP for P, and versus K+P for MSP. Mean matched-cell differences were summarized by 10,000 paired bootstrap resamples. Two-sided paired Wilcoxon signed-rank tests were adjusted across the 12 metric-by-contrast rows by the Benjamini--Hochberg procedure \citep{Wilcoxon1945,Benjamini1995}. Same-dimension PCA/SVD controls addressed dimensionality effects, and block-weight audits tested reasonable reweightings of the composition and property blocks.

The mixed-feature question was also tested through empirical mutual-information (MI) audits. For a perturbation label $Y$ and representation-derived distance features $X$, an equal-frequency discretized screen was retained as a supplementary descriptive check. The primary analysis used a Ross/Kraskov--St\"ogbauer--Grassberger (KSG)-style nearest-neighbour estimator with $k=5$ on low-dimensional continuous summaries and a discrete local-versus-nuisance label \citep{Kraskov2004,Ross2014}. It evaluated $d_K$, $d_K+d_P$, $d_K+d_M$ and $d_K+d_P+d_M$ across 12 local-mutation cells, with 100 label permutations and 100 stratified subsamples per cell. These estimates are estimator-dependent empirical summaries rather than absolute physical information quantities. The seven-group contribution audit was paired with a relation audit based on rank, canonical-correlation analysis (CCA) and row-wise P/MSP alignment. Together, these analyses test conditional empirical roles; they do not imply that P and MSP measure unrelated physical quantities.

We added a dimension-matched high-k compression audit to avoid comparing CK4P-MSP only with short contiguous k-mer baselines. Canonical $k=15$ signals were compressed to the 222-feature CK4P-MSP budget using hashing-trick counts and sparse random projection. These vector controls were evaluated on the same 69, 75, 100 and 150 bp perturbation grid using paired cosine similarity, L2 drift, nearest-clean retrieval and shallow readout. MinHash was assessed separately in its native estimated-versus-exact Jaccard geometry and was not treated as an L2 vector or linear-readout feature. This audit tests whether compact stability persists against high-specificity k-mer information under a comparable vector budget. Finally, a feature-extraction runtime audit compared CK4P-MSP with CK4, CK4+P and the compact high-k controls. Runtime was treated as an engineering boundary, not as a primary performance claim.

\subsection{Reproducibility}

The public implementation and contract-v2 result namespace are maintained in the release repository. The manuscript-facing API is \path{methods/ck4p_msp.py}; the contribution, high-k, KSG, MSP-sensitivity and redundancy/runtime audits are executed from \path{experiments/audits/} and mapped to their result tables in \path{docs/contract_v2_evidence_map.md}. Historical scripts and outputs remain available for provenance but are not aliases for the public block-normalized CK4P-MSP method. The supported scope is compact and controlled: biochemical and position-aware summaries are evaluated as block-decomposable representation-level audit coordinates in short-read settings.

\FloatBarrier
\section{Results}

\subsection{Empirical mixed-feature information audit under local perturbation}

We first tested whether the added property channels showed perturbation-associated signal rather than only changing the scale of a mixed feature space. The contract-v2 KSG-style audit was applied to low-dimensional blockwise distance summaries for local change versus matched nuisance perturbation. Across 12 local-mutation cells, the mean estimate was 0.889 bits for $d_K+d_P+d_M$ and 0.752 bits for $d_K$ alone, an average incremental estimate of 0.138 bits. All primary summaries exceeded their label-permutation reference under the tested audit. This result is estimator-dependent empirical evidence: under the tested perturbation grid and estimator, P/MSP distance summaries improved local-versus-noise separability beyond the CK4 distance summary.

This audit does not turn the mixed representation into an information-theoretic theorem. It provides a complementary check on the ablation results: the property blocks contribute empirically detectable perturbation signal under the tested local-change regime, while the block-normalized metric remains a standardized diagnostic drift rather than a physical distance with commensurate units.

Estimator sensitivity and permutation references support this result (Supplementary Figure S2).

\subsection{Global and positional property layers have related but distinct roles}

We then evaluated all seven non-empty K/P/MSP combinations over 69, 75, 100 and 150 bp reads. Pure property summaries had the lowest numerical drift (P, 0.025; MSP, 0.027; P+MSP, 0.027), but they did not preserve the same nearest-clean behavior as K-containing representations: retrieval was 0.295 for P and 0.942 for P+MSP, compared with 1.000 for CK4P-MSP. The composition block therefore traded some numerical stability for composition-linked retrieval. Relative to P+MSP, adding K increased retrieval by 0.058 (95\% bootstrap CI 0.036--0.081; BH-adjusted $q=7.3\times10^{-4}$) while increasing, rather than decreasing, drift.

Conditional effects differed by audit axis. Adding P to CK4+MSP reduced mean drift from 0.156 to 0.129, a matched-cell improvement of 0.027 (95\% CI 0.023--0.031; $q=2.4\times10^{-7}$), but did not materially change grouped delta-readout (0.9780 versus 0.9786; $q=0.525$). Adding MSP to CK4+P produced a similar drift reduction of 0.027 (95\% CI 0.023--0.032; $q=2.4\times10^{-7}$) and increased grouped local-change macro-F1 from 0.904 to 0.979 (mean difference 0.074, 95\% CI 0.035--0.117; $q=7.3\times10^{-4}$). Adding K to P+MSP increased macro-F1 by only 0.008 and the interval included zero. These results assign P primarily to global stability, MSP to local-change readability and K to composition-linked retrieval under the tested grid; they do not establish statistical independence or universal necessity.

The redundancy audit also cautions against overinterpreting P and MSP as separable sources. Across the tested lengths, the first canonical correlation between P and MSP was near one and their first principal components were strongly aligned, as expected because both blocks are derived from the same biochemical property family. However, MSP had higher numerical rank than P, and row-wise P/MSP alignment after compression was incomplete. We therefore interpret the blocks as related layers at different scales, not as interchangeable evidence sources or as measurements of unrelated physical quantities.

\begin{figure*}[!tbp]
\centering
\includegraphics[width=0.92\textwidth]{figures/main/nature_fig2_compact_stability.pdf}
\caption{Seven-group K/P/MSP ablation across 24 matched length-by-perturbation cells and 12 grouped local-mutation cells. (A) Mean paired L2 drift is lowest for property-only summaries; CK4P-MSP is more stable than CK4 while retaining the K block. (B) MSP-containing representations preserve the strongest grouped local-versus-nuisance delta-readout. (C) K-containing representations retain near-perfect nearest-clean retrieval, whereas P-only and P+MSP lose composition-linked retrieval. Points show cell means; exact paired bootstrap intervals, Wilcoxon tests and Benjamini--Hochberg-adjusted q values are reported in Supplementary Table S2.}
\label{fig:figure2}
{\small\noindent\textbf{Alt text:} Three aligned dot plots compare all seven K/P/MSP block combinations by global paired L2 drift, grouped local-change macro-F1 and nearest-clean retrieval. Property-only summaries have the lowest drift, MSP-containing summaries have the strongest local-change readout, and K-containing summaries have the strongest retrieval.\par}
\end{figure*}


The P/MSP relation and runtime audit is shown in Supplementary Figure S8.

\subsection{MSP short-bin reliability remains usable at 69-75 bp}

We next tested whether MSP becomes noise-dominated when ultra-short reads are divided into finer relative-position bins. The short-bin audit did not show an immediate loss of reliability at 69 or 75 bp. At the pre-specified $\gamma=1$ setting, the full 2+3+4+6 binset had paired cosine of approximately 0.985 and L2 drift of 0.153 at 69 bp, and paired cosine of approximately 0.987 and L2 drift of 0.140 at 75 bp. Retrieval remained effectively saturated.

The grouped local delta-readout audit did not show an immediate short-bin failure. At 69 bp and $\gamma=1$, mean macro-F1 increased from 0.876 with 2 bins to 0.907 with 2+3 bins, 0.951 with 2+3+4 bins and 0.959 with 2+3+4+6 bins. The same monotonic pattern persisted at 100 and 150 bp. Increasing $\gamma$ to 2 further reduced global drift, so the default $\gamma=1$ should be understood as a fixed balanced setting rather than an empirically optimal stability weight. The combined interpretation is restrained but favorable: static MSP did not become noise-dominated in the tested ultra-short-read regime, yet it remains a coarse positional summary rather than an optimized variance-aware weighting model.

MSP binset and weight sensitivity are shown in Supplementary Figure S6.

\subsection{Error-aware ART perturbation audit}

We next asked whether the external simulator evidence depended on treating ART as one pooled error source. In the quality-stratified ART audit, representative compact controls remained readable across low-, mid- and high-quality bins. The property-aware compact representation preserved high paired cosine and bounded L2 drift across quality strata, and the spaced-property comparator retained retrieval and readout consistency.

These results strengthen, but do not overstate, the ART evidence. They show that the main pattern is not restricted to an undifferentiated simulator summary. They do not claim that ART quality strata capture the full physical error distribution of a clinical sequencing run.

Quality-stratified ART results are shown in Supplementary Figure S3.

\subsection{Compact biochemical summaries reduce perturbation drift}

The main compact-representation audit tested whether property-aware summaries remained closer to their clean counterparts than composition-only or compressed high-k baselines. Across the WGS-derived perturbation grid, CK4P-MSP retained paired cosine of 0.989 and mean L2 drift of 0.129 at 222 features, compared with 0.969 and 0.219 for CK4 (136 features) and 0.946 and 0.293 for CK5 (512 features). CK4+P was more globally stable than CK4P-MSP (paired cosine 0.996; drift 0.078), which is consistent with its role as a global property stability summary rather than a positional representation.

Dimension-matched high-k vector controls further addressed whether a conventional high-specificity k-mer representation could show the same behavior under the CK4P-MSP feature budget. At approximately 222 features, CK4P-MSP retained paired cosine 0.989 and L2 drift 0.129, compared with 0.821 and 0.524 for hashed $k=15$ counts and 0.868 and 0.447 for sparse random projection of $k=15$ counts. Shallow-readout differences were modest and did not favour CK4P-MSP over all high-k controls, so the result supports a compact perturbation-stability advantage rather than a predictive-performance claim. MinHash was assessed separately through its native estimated-versus-exact Jaccard agreement (mean estimated Jaccard 0.711, exact Jaccard 0.711, mean absolute error 0.021); it was not used as an L2-vector baseline. Paired Wilcoxon tests over matched analysis cells supported the primary stability contrasts after multiple-testing correction. Together, these controls show that the stability trend is not explained solely by feature count or by an arbitrary global feature projection.

\begin{figure*}[!tbp]
\centering
\includegraphics[width=0.92\textwidth]{figures/main/nature_fig3_ck4p_msp_tradeoff.pdf}
\caption{Dimension-matched compact stability audit over the controlled WGS perturbation grid. Panels compare mean paired cosine similarity and standardized L2 drift across compact composition/property representations and 222-dimensional hashed or sparse-random-projection $k=15$ vectors, and place stability against feature dimension. CK4P-MSP is more stable than CK4, CK5 and both compressed high-k vector controls; CK4+P remains the most stable compact global-summary representation. MinHash is excluded from this L2-vector comparison and evaluated in native Jaccard geometry.}
\label{fig:figure3}
{\small\noindent\textbf{Alt text:} Bar plots and a scatter plot comparing compact k-mer, property-aware and dimension-matched high-k vector representations by paired similarity, standardized diagnostic drift and feature dimension.\par}
\end{figure*}


The same-dimension and mixed-metric baseline controls are summarized in Supplementary Figure S1.

Estimator and compressed high-k sensitivity results are provided in Supplementary Figure S7.

\subsection{CK4P-MSP balances compactness, stability and representation readability}

We next evaluated CK4P-MSP as a compact trade-off rather than as a single-metric winner. The seven-group result showed why the complete representation contains all three blocks: property-only variants were numerically stable but weakened nearest-clean retrieval, P improved global stability after conditioning on K+MSP, and MSP improved local-change readout after conditioning on K+P. CK4P-MSP therefore preserves a compact canonical local-composition backbone while exposing global biochemical and coarse positional audit channels.

The resulting trade-off supported the intended role. CK4P-MSP remained compact, reduced perturbation drift relative to CK4 and CK5 composition baselines, retained near-perfect nearest-clean retrieval and preserved accessible local-change signal. Its role is to provide a block-decomposable audit coordinate system between composition-only compact features and high-dimensional full-position matrices, rather than to maximize predictive accuracy against learned embeddings.

\begin{table*}[!tbp]
\centering
\scriptsize
\caption{Compact main-method metrics across the controlled WGS and grouped local-change audits. Paired cosine and L2 drift are averaged over the WGS perturbation grid; logistic macro-F1 is a shallow accessibility probe rather than production-classifier performance. Feature count reports the fixed representation dimension.}
\label{tab:table3}
\begin{tabular}{@{}p{0.13\textwidth}p{0.22\textwidth}p{0.21\textwidth}p{0.14\textwidth}p{0.13\textwidth}@{}}
\toprule
\textbf{representation} & \textbf{paired cosine} & \textbf{L2 drift} & \textbf{readout macro-F1} & \textbf{features} \\
\midrule
CK4 & 0.935 & 0.316 & 0.378 & 136 \\
CK4+P & 0.990 & 0.122 & 0.380 & 147 \\
CK4P-MSP & 0.991 & 0.116 & 0.401 & 222 \\
CK5 & 0.882 & 0.433 & 0.373 & 512 \\
\bottomrule
\end{tabular}

\end{table*}


\subsection{External ART and CAMI probes separate stability from task-limited readability}

ART and CAMI were used to test whether the representation family showed consistent behaviour outside the handcrafted perturbation grid, while keeping stability and readout as separate questions. ART provided simulator-derived paired-error context, CAMI\_TOY\_low provided a lightweight external coarse target-versus-background readout probe, and CAMI II marine provided an anonymous external stability probe. These legacy external probes remain useful contextual checks, but they were generated before the public block-normalized method contract was introduced. They therefore do not supply the primary quantitative CK4P-MSP claims in this revision; those claims are restricted to the contract-v2 WGS perturbation, ablation and grouped local-readout analyses.

\begin{figure*}[!tbp]
\centering
\includegraphics[width=0.92\textwidth]{figures/main/nature_fig4_external_probes.pdf}
\caption{External probes separate stability from task-limited readability. (A) Mean paired cosine under ART Illumina-like simulation tests simulator-derived stability. (B) CAMI\_TOY\_low target-versus-background macro-F1 at 69 and 100 bp tests coarse external readability; these historical readouts predate the public block-normalized method contract. (C) Mean L2 drift on anonymous CAMI II marine reads tests external perturbation stability without reconstructed read-level taxonomic labels. Lower drift is preferable in panel C; panels B and C are not taxonomic validation of contract-v2 CK4P-MSP.}
\label{fig:figure4}
{\small\noindent\textbf{Alt text:} Three-panel comparison of ART stability, CAMI coarse target-background readout and CAMI II anonymous-read stability.\par}
\end{figure*}


The quality-stratified ART audit and the CAMI II marine subset remain supplementary external checks. The former stratifies a simulator-derived error layer by read quality; the latter applies paired perturbations to anonymous public marine reads without reconstructing read-level taxonomic labels. Marine composition can differ materially from pathogen-rich or low-biomass mNGS sources. These probes therefore provide context for the breadth of the perturbation setting, not external taxonomic validation or a contract-v2 method comparison.

Together, the external resources motivate a future public-API rerun with a larger labelled benchmark. In the present revision, ART supports simulator context, CAMI\_TOY\_low supports task-limited external readability, and CAMI II supports anonymous-read stability only; the primary quantitative evidence remains the contract-v2 controlled WGS and grouped local-mutation analyses.

The complete CAMI II marine length-by-perturbation matrix is provided in Supplementary Figure S10.

\subsection{Full-position matrices define a positional upper bound rather than a deployment default}

Full-position encodings were used to estimate what is gained when fine-grained layout information is retained. In controlled position and order tasks, full-position token or property matrices improved positional readability compared with compact summaries. This was expected, because these matrices preserve per-position structure that compact summaries intentionally compress.

\begin{figure*}[!tbp]
\centering
\includegraphics[width=0.92\textwidth]{figures/main/nature_fig5_full_position_upper_bound.pdf}
\caption{Full-position diagnostic upper bound for positional information. (A) Macro-F1 on controlled positional readout tasks compares compact summaries with representations that retain per-position token or property channels. (B) Readout is plotted against feature dimension to show the cost of preserving near-explicit positional information. CK4P-MSP provides a compact coarse summary and is not the positional-readout upper bound.}
\label{fig:figure5}
{\small\noindent\textbf{Alt text:} Bar and scatter plots showing that full-position matrices add positional readout value at substantially higher feature dimension.\par}
\end{figure*}


The result is best interpreted as an upper-bound diagnostic. Full-position matrices clarify the value of positional information, but their high dimensionality and task specificity make them unsuitable as the default representation for compact short-read auditing. CK4P-MSP therefore occupies a different role: it does not match the full positional upper bound, but it retains a coarse positional property summary at a much smaller feature cost.

\subsection{Local mutation analysis separates robustness from selective sensitivity}

A stable audit representation should remain responsive to structured changes, but the relevant notion of responsiveness depends on the metric. We therefore tested whether local biochemical change remained distinguishable from matched nuisance perturbation. In a distance-ratio analysis, CK4P-MSP did not exceed CK4: its local-mutation/random-noise L2 ratio was 0.752 compared with 0.750 for CK4. Full-position and property-channel representations were stronger distance-amplification probes. CK4P-MSP's positive result was different: it retained high grouped delta-readout at compact dimensionality, with local-versus-noise macro-F1 of 0.979. The seven-group ablation assigned most of this local-readout increment to MSP, while P contributed primarily to global stability and K to nearest-clean retrieval.

\begin{figure*}[!tbp]
\centering
\includegraphics[width=0.92\textwidth]{figures/main/nature_fig6_local_mutation_sensitivity.pdf}
\caption{Local-change readout and distance-ratio boundary. (A) Grouped local-versus-nuisance macro-F1 across the compact K/P/MSP combinations shows that MSP accounts for most of the local-change readout increment. All derivatives of one template remained in the same cross-validation fold. (B) CK4P-MSP does not exceed CK4 on the raw local-to-nuisance L2 ratio, marking a boundary rather than a claimed mutation-distance advantage.}
\label{fig:figure6}
{\small\noindent\textbf{Alt text:} Bar plots showing MSP-driven grouped local-change readability and the separate compact-representation distance-ratio boundary.\par}
\end{figure*}


\FloatBarrier

This result supports a division of labor. Paired cosine and L2 drift quantify nuisance stability, distance ratios identify upper-bound local-change probes, and delta-readout asks whether structured local changes remain accessible to a shallow model. CK4P-MSP should therefore not be described as a distance-amplifying mutation detector. Its value in this analysis is that it remains stable under nuisance perturbation while preserving decomposable information for local-change readout at much lower dimensionality than full-position matrices.

\begin{table*}[!tbp]
\centering
\scriptsize
\caption{Seven-group block ablation. Global drift and nearest-clean retrieval are averaged over 24 matched length-by-perturbation cells. Grouped delta-readout macro-F1 is averaged over 12 length-by-local-mode cells using template-grouped fivefold cross-validation. Lower drift and higher retrieval/readout indicate stronger performance on their respective audit axes.}
\label{tab:table4}
\begin{tabular}{@{}p{0.126\textwidth}p{0.126\textwidth}p{0.126\textwidth}p{0.126\textwidth}p{0.126\textwidth}p{0.126\textwidth}p{0.126\textwidth}@{}}
\toprule
\textbf{representation} & \textbf{noise L2} & \textbf{local L2} & \textbf{local -- noise L2} & \textbf{sensitivity ratio} & \textbf{delta-readout macro-F1} & \textbf{features} \\
\midrule
CK4 & 0.355 & 0.266 & -0.089 & 0.750 & 0.891 & 136 \\
CK4+P & 0.115 & 0.086 & -0.029 & 0.752 & 0.900 & 147 \\
CK4P-MSP & 0.110 & 0.083 & -0.027 & 0.754 & 0.977 & 222 \\
CSP & 0.116 & 0.109 & -0.006 & 0.947 & 0.696 & 147 \\
one-hot & 0.239 & 0.239 & 0.000 & 1.000 & 0.956 & 500 \\
P-channels & 0.086 & 0.107 & 0.020 & 1.255 & 0.983 & 500 \\
one-hot+P & 0.188 & 0.207 & 0.019 & 1.101 & 0.973 & 1000 \\
position k-mer P & 0.167 & 0.185 & 0.017 & 1.125 & 0.972 & 1536 \\
\bottomrule
\end{tabular}

\end{table*}


The mutation-fraction sweep is shown in Supplementary Figure S4.

\subsection{Boundary analyses constrain the interpretation}

The boundary probes were included to prevent overextension of the representation claim. Spaced-seed transfer was useful as a mechanism comparator, but matching-oriented seed priors did not automatically become a universal dense-feature advantage. Context visibility depended on whether the relevant relation was present within the read. ARG/SNP interpretation required identity, annotation and allele-level evidence beyond what compact property summaries can provide.

These negative and boundary results are part of the manuscript's logic. They keep the claim focused on representation diagnostics and representation-level perturbation auditing, rather than allowing the compact feature space to be misread as a full biological interpretation system.

\begin{table*}[!tbp]
\centering
\scriptsize
\caption{Boundary and mechanism summary.}
\label{tab:table5}
\begin{tabular}{@{}p{0.22\textwidth}p{0.36\textwidth}p{0.34\textwidth}@{}}
\toprule
\textbf{boundary} & \textbf{main observation} & \textbf{interpretation} \\
\midrule
spaced-seed transfer & The best tested 4-position dense stability pattern was contiguous 0-1-2-3 across the 69/75 bp sanity grid. & spaced seeds remain useful matching priors, but transfer is regime-specific \\
context visibility & motif-pair visibility thresholds depended on placement, emerging at 130, 140, 148 or 155 bp. & a representation cannot encode context absent from the read \\
ARG/SNP readout & stable compact features did not by themselves establish allele, resistance-SNP or functional equivalence. & identity/database evidence remains necessary for final biological calls \\
\bottomrule
\end{tabular}

\end{table*}
